% =====================================================================
% DECEPTION COMPETITION — METHODOLOGY (Option-R design), paper-ready.
% Supersedes the audit-design deception_section.tex. Pairs with the results in
% deception_results_R.tex. Style mirrors auction_section_rewrite.tex.
%
% Key R changes vs the audit design: FULL information (no partial-info sweep, no
% per-attribute audit / catch event / population mean); score = (w.c) * trust
% (no 0.5 floor); PERSISTENT compounding credibility updated for every agent
% every round; no booking threshold. Mimic input is 14-dim with a behavioral
% history block; inference adds calibrated residual noise.
% =====================================================================

\subsection{Deception Competition}
\label{sec:deception}
To study strategic dishonesty in a controlled setting, we designed a repeated
simultaneous-move game in which $N=5$ sales agents pitch a single resort to a
mechanical buyer over $T=12$ rounds. Each round the resort is described by $K=5$
quality attributes whose true values form a vector $t^{(r)}\in[0,1]^{K}$,
quantized to two decimal places and \emph{re-sampled each round}. Every agent
observes the full truth $t^{(r)}$ (full information), then simultaneously emits a
\emph{claim} vector $c_i^{(r)}\in[0,1]^{K}$. Because every agent sees the truth,
a claim above it is a deliberate over-statement rather than a guess, and the
game is purely one of how far each agent will distort what it knows given that
distortion erodes its standing with the buyer. The buyer holds fixed preference
weights $w\in\Delta^{K}$ (default $w=(0.25,0.20,0.20,0.20,0.15)$) and is
non-strategic.

\subsubsection{Scoring and booking}
The buyer scores each agent by its trust-weighted appeal, where the standing
trust $\mathrm{tr}_i^{(r)}\in[0,1]$ (initialized to $1$) directly scales the
appeal:
\begin{equation}
s_i^{(r)} = \bigl(w \cdot c_i^{(r)}\bigr)\,\mathrm{tr}_i^{(r)}.
\label{eq:dec-r-score}
\end{equation}
Trust enters as a bare multiplier (no floor): a fully discredited agent
($\mathrm{tr}=0$) scores $0$ and cannot win, which is what makes sustained
dishonesty self-defeating. The buyer books the highest-scoring agent each round,
\begin{equation}
\mathcal{W}^{(r)} = \operatorname*{arg\,max}_{i}\, s_i^{(r)},
\label{eq:dec-r-winner}
\end{equation}
set-valued under ties, with no minimum threshold. The per-round reward pays
winners an equal split and everyone else nothing,
\begin{equation}
R_i^{(r)} =
\begin{cases}
1/\lvert \mathcal{W}^{(r)} \rvert, & i \in \mathcal{W}^{(r)},\\[2pt]
0, & \text{otherwise,}
\end{cases}
\label{eq:dec-r-reward}
\end{equation}
so an agent's episode return $R_i=\sum_r R_i^{(r)}$ is simply its number of
(tie-shared) bookings, bounded by $T$. The score magnitude never enters the
reward directly; trust matters only through who wins.

\subsubsection{Credibility dynamics}
The single cost of lying is a \emph{persistent} loss of credibility. After each
round the buyer compares every agent's claim to the now-revealed truth and
computes its weighted \emph{over-promise} (disappointment),
\begin{equation}
d_i^{(r)} = \max\!\Bigl(0,\ \textstyle\sum_{a} w_a\,\bigl(c_{i,a}^{(r)} - t_a^{(r)}\bigr)\Bigr),
\label{eq:dec-r-disc}
\end{equation}
which ignores under-claiming (the buyer is never disappointed by a too-modest
pitch). Credibility then drifts up with honesty and erodes with over-promise,
for \emph{every} agent on \emph{every} round, clipped to $[0,1]$:
\begin{equation}
\mathrm{tr}_i^{(r+1)} = \mathrm{clip}\!\bigl(\mathrm{tr}_i^{(r)} + \delta - \lambda\,d_i^{(r)},\ 0,\ 1\bigr),
\label{eq:dec-r-trust}
\end{equation}
with recovery $\delta=0.1$ and penalty rate $\lambda=1.5$. An honest round adds
$\delta$; an over-claiming round subtracts $\lambda d$. Because the update is
applied to all agents every round (not only winners), \emph{sustained}
over-claiming compounds credibility toward $0$, whereas a single lie is
recoverable over a few honest rounds. The constants are set so that the
math-tier reference ladder (below) forms cleanly---naive constant inflation is
punished below honest play while \emph{selective} inflation still pays---making
the game a test of \emph{calibrated} deception rather than of inflation volume.

\subsubsection{Behavioral metrics}
\label{sec:dec-r-metrics}
We summarize an agent's revealed strategy with the \emph{lie rate}---the fraction
of claimed attributes that exceed the truth,
\begin{equation}
\mathrm{LieRate}_i =
\frac{\sum_r \sum_{a} \mathbb{1}\!\left[c_{i,a}^{(r)} > t_a^{(r)}\right]}{T\,K},
\label{eq:dec-r-lierate}
\end{equation}
the win rate $\mathrm{WinRate}_i = \tfrac{1}{T}\sum_r \mathbb{1}[i \in \mathcal{W}^{(r)}]$,
the per-round over-promise $d_i^{(r)}$ of Eq.~\eqref{eq:dec-r-disc}, and the
final credibility $\mathrm{tr}_i^{(T)}$. To probe \emph{calibration} we also bin
$d_i^{(r)}$ by the trust $\mathrm{tr}_i^{(r)}$ an agent carries into the round: a
strategic agent reduces over-promise as credibility falls, so the trend of
$d$ against incoming trust distinguishes calibrated from indiscriminate
deception.

\subsubsection{Math-tier reference strategies}
\label{sec:dec-r-tiers}
We define a sophistication-ordered ladder of hand-coded policies, all acting on
the full truth $t^{(r)}$, own trust, and round index. Writing
$a^{\star}=\operatorname*{arg\,max}_a w_a(1-t_a)$ for the highest headroom-weighted
attribute and $\phi=r/(T-1)$ for the round fraction:
\begin{itemize}
\item \textbf{T1 Honest:} $c=t$. Never over-claims; trust stays $1$.
\item \textbf{T2 Naive inflate:} $c=\mathrm{clip}(t+0.2)$. Over-claims every
attribute every round; the broken control.
\item \textbf{T3 Trust-conditional:} inflate $a^{\star}$ to $1$ \emph{iff}
$\mathrm{tr}>0.6$, else play honestly---lie only while credibility can afford it,
and stop to recover otherwise.
\item \textbf{T4 Backloaded:} honest until the back half, then ramp the inflation
of $a^{\star}$ by $\min(1,\max(0,(\phi-0.5)/0.5))$---time the lie to the end-game,
when credibility damage has few remaining rounds to compound.
\end{itemize}
T3 and T4 are the two ``calibrated'' tiers; T1 is the safe baseline and T2 the
naive corner.

\subsubsection{Behavioral mimics}
\label{sec:dec-r-mimics}
As in the auction study, we distill each LLM into a cheap surrogate and use it as
a drop-in replacement at scale. At agent $i$'s decision point we encode the
leakage-safe public context as a normalized vector
$\mathbf{x}_i^{(r)}\in\mathbb{R}^{14}$ concatenating the truth ($K=5$), the
agent's own trust, the four opponents' trusts, the round fraction $\phi$, and a
compact behavioral history block (the agent's previous-round over-promise, a
won-last-round indicator, and its win rate so far). The history block is what
lets the surrogate reproduce timing-dependent play such as back-loading. The
vector is z-scored and passed through a shared two-layer trunk (width $32$,
ReLU, dropout $0.2$) feeding two heads,
\begin{align}
p_a &= \sigma\!\bigl((W_\ell h + b_\ell)_a\bigr), &
\hat{c}_a &= \sigma\!\bigl((W_c h + b_c)_a\bigr),
\label{eq:dec-r-heads}
\end{align}
the \emph{lie head} $p_a$ (probability of over-claiming attribute $a$) and the
\emph{claim head} $\hat{c}_a$ (regressed claim value). They are trained with
binary cross-entropy on the lie labels and mean-squared error on the claim values
restricted to attributes that were lied on,
\begin{equation}
\begin{aligned}
\mathcal{L} &= \mathrm{BCE}(p, y^{\mathrm{lie}})
 + \lambda_c\,\frac{\sum_a y^{\mathrm{lie}}_a\,(\hat{c}_a - y^{\mathrm{claim}}_a)^2}{\sum_a y^{\mathrm{lie}}_a},
\end{aligned}
\label{eq:dec-r-loss}
\end{equation}
with $\lambda_c=1$. At deployment the mimic draws a lie decision
$\ell_a\sim\mathrm{Bernoulli}(p_a)$; on attributes it does not lie about it snaps
to the truth, and otherwise emits the regressed value perturbed by calibrated
Gaussian noise,
\begin{equation}
c_a =
\begin{cases}
t_a, & \neg\ell_a,\\[2pt]
\mathrm{clip}_{[0,1]}\!\bigl(\hat{c}_a + \varepsilon_a\bigr),\ \varepsilon_a\sim\mathcal{N}\!\bigl(0,(\Theta\hat{\sigma}_a)^2\bigr), & \ell_a,
\end{cases}
\label{eq:dec-r-infer}
\end{equation}
where $\hat{\sigma}_a$ is the claim head's per-attribute training residual
standard deviation, saved with the checkpoint. This noise is essential: without
it the regressor collapses to the conditional mean and the mimic's claims are
under-dispersed relative to the real LLM's.

\subsubsection{Fidelity validation and the transfer test}
\label{sec:dec-r-fidelity}
We validate the substitution distributionally: we bin each attribute's claims and
compare the real and mimic histograms $P,Q$ by total-variation distance,
\begin{equation}
\mathrm{TVD}(P,Q) = \tfrac12\sum_b \lvert P_b - Q_b\rvert,
\label{eq:dec-r-tvd}
\end{equation}
used as the acceptance gate because it is bounded and sample-size independent (a
$\chi^2$ $p$-value rejects negligible differences at our sample sizes). Episodes
are the independent replicate and we report $95\%$ bootstrap confidence intervals
over episodes throughout. Crucially, because distributional fidelity need not
imply that strategy \emph{rankings} measured against the mimics hold against the
real models, we additionally run a \emph{transfer test}: each reference strategy
is played against a fixed field of behavioral mimics and, separately, against the
real LLMs, and the two strategy orderings are compared. Agreement licenses the
mimics as a comparison tool; disagreement is itself reported (Sec.~\ref{sec:dec-r-transfer}).
